\chapter{WP1: Sun-Earth and Sun-Mars Lagrange Points}
\label{chapter:wp1}

\section{Task 1.1}

\subsection{Task 1.1a}

\begin{table}[H]
  \centering
  \caption{Dimensionless coordinates $(x_i, y_i, z_i)$ of the Sun-Earth and Sun-Mars collinear Lagrange points in their respective synodic reference frames.}
  \label{tab:lagrange_coords}
  \begin{tabular}{@{} l S[table-format=1.10] c c @{}}
    \toprule
    \textbf{Lagrange Point} & {$x_i$} & {$y_i$} & {$z_i$} \\
    \midrule
    Sun-Earth - $L_2$ & 1.0100904358 & 0 & 0 \\
    Sun-Mars  - $L_2$ & 1.0047631067 & 0 & 0 \\
    \bottomrule
  \end{tabular}
\end{table}

\subsection{Task 1.1b}

The Lagrange points are verified using \textcite{wakker2015fundamentals}'s characteristic distance parameter $\gamma_2 = x_{L_2} - (1-\mu)$.
For Sun-Mars L2, \citeauthor{wakker2015fundamentals} tabulates $\gamma_2 = 4.763 \times 10^{-3}$; our solver returns $\gamma_2 = 4.763353 \times 10^{-3}$, agreeing to the sixth decimal place (difference: $3.53 \times 10^{-7}$).
The Sun-Earth case is verified analogously.

(46 words)

\section{Task 1.2} \label{sec:1.2}

At L2, three acceleration components act on the spacecraft in the \ac{cr3bp} rotating frame: the Sun's gravity ($-x$), the planet's gravity ($-x$), and the centrifugal acceleration ($+x$) \parencite{wakker2015fundamentals,koon2000dynamical}. 
Equilibrium requires the outward centrifugal term to balance the combined inward gravitational pulls. 
Earth's mass ratio is roughly one order of magnitude larger than Mars's, so its gravitational reach extends farther; the planet--L2 distance scales as the Hill radius, $\gamma \approx \left(\frac{\mu}{3}\right)^{1/3}$ \parencite{szebehely1967theory}. % this cite comes through Koon et al., equation 2.5.3
Hence
\begin{equation}
  \gamma_E \approx 0.01009 > \gamma_M \approx 0.00476 \Rightarrow \frac{\gamma_E}{\gamma_M} \approx \left( \frac{\mu_E}{\mu_M} \right)^{\frac{1}{3}} \approx 2.12.
\end{equation} 

(67 words) % maybe too short?